\documentclass[12pt, a4paper]{article}

% ============================================================
% Pakete
% ============================================================
\usepackage[T1]{fontenc}
\usepackage[utf8]{inputenc}
\usepackage[ngerman]{babel}
\usepackage{geometry}
\geometry{margin=2.5cm}
\setlength{\headheight}{15pt}

\usepackage{lmodern}
\usepackage{microtype}
\usepackage{parskip}
\usepackage{graphicx}
\usepackage{xcolor}
\usepackage{listings}
\usepackage{booktabs}
\usepackage{longtable}
\usepackage{array}
\usepackage{tabularx}
\usepackage{hyperref}
\usepackage{fancyhdr}
\usepackage{titlesec}
\usepackage{enumitem}
\usepackage{mdframed}
\usepackage{tcolorbox}
\usepackage{amsmath}
\usepackage{amssymb}
\usepackage{float}
\usepackage{caption}
\usepackage{subcaption}
\usepackage{tikz}
\usetikzlibrary{shapes.geometric, arrows.meta, positioning, fit, backgrounds, calc}

% ============================================================
% Farben
% ============================================================
\definecolor{codebg}{RGB}{245,245,245}
\definecolor{codeframe}{RGB}{200,200,200}
\definecolor{keywordcolor}{RGB}{0,0,180}
\definecolor{commentcolor}{RGB}{100,130,80}
\definecolor{stringcolor}{RGB}{180,0,0}
\definecolor{accentblue}{RGB}{0,90,160}
\definecolor{accentgray}{RGB}{90,90,90}
\definecolor{boxbg}{RGB}{235,245,255}
\definecolor{boxframe}{RGB}{0,90,160}

% ============================================================
% Listings-Stil
% ============================================================
% UTF-8-Unterstuetzung fuer listings (Umlaute + Sonderzeichen)
\lstset{
  literate=%
    {ä}{{\"{a}}}1 {ö}{{\"{o}}}1 {ü}{{\"{u}}}1
    {Ä}{{\"{A}}}1 {Ö}{{\"{O}}}1 {Ü}{{\"{U}}}1
    {ß}{{\ss{}}}1
    {→}{{$\rightarrow$}}1
    {←}{{$\leftarrow$}}1
    {…}{{...}}1
    {│}{{|}}1
    {├}{{+}}1
    {└}{{+}}1
    {─}{{-}}1
    {✓}{{$\checkmark$}}1
    {❌}{{X}}1
}

\lstdefinestyle{pythonstyle}{
  language=Python,
  backgroundcolor=\color{codebg},
  frame=single,
  rulecolor=\color{codeframe},
  basicstyle=\ttfamily\small,
  keywordstyle=\color{keywordcolor}\bfseries,
  commentstyle=\color{commentcolor}\itshape,
  stringstyle=\color{stringcolor},
  showstringspaces=false,
  breaklines=true,
  breakatwhitespace=true,
  numbers=left,
  numberstyle=\tiny\color{accentgray},
  numbersep=8pt,
  tabsize=4,
  captionpos=b,
  xleftmargin=1.5em,
}

\lstdefinestyle{yamlstyle}{
  language=,
  backgroundcolor=\color{codebg},
  frame=single,
  rulecolor=\color{codeframe},
  basicstyle=\ttfamily\small,
  keywordstyle=\color{keywordcolor}\bfseries,
  commentstyle=\color{commentcolor}\itshape,
  stringstyle=\color{stringcolor},
  showstringspaces=false,
  breaklines=true,
  numbers=left,
  numberstyle=\tiny\color{accentgray},
  numbersep=8pt,
  tabsize=2,
  captionpos=b,
  xleftmargin=1.5em,
  morecomment=[l]{\#},
  morestring=[b]",
  morestring=[b]',
}

% ============================================================
% tcolorbox-Stile
% ============================================================
\tcbuselibrary{skins, breakable}

\newtcolorbox{infobox}[1][]{
  colback=boxbg,
  colframe=boxframe,
  fonttitle=\bfseries,
  title={#1},
  breakable,
}

\newtcolorbox{definitionbox}[1][]{
  colback=white,
  colframe=accentgray,
  fonttitle=\bfseries\color{white},
  coltitle=white,
  colbacktitle=accentgray,
  title={#1},
  breakable,
}

% ============================================================
% Seitenlayout
% ============================================================
\pagestyle{fancy}
\fancyhf{}
\rhead{\textcolor{accentgray}{\small KI-Copilot für Verzahnungswissen}}
\lhead{\textcolor{accentgray}{\small Systemdokumentation}}
\cfoot{\thepage}
\renewcommand{\headrulewidth}{0.4pt}

\hypersetup{
  colorlinks=true,
  linkcolor=accentblue,
  urlcolor=accentblue,
  citecolor=accentblue,
  pdftitle={KI-Copilot für Verzahnungswissen -- Systemdokumentation},
  pdfauthor={Automatisch generiert},
}

% ============================================================
% Titelformatierung
% ============================================================
\titleformat{\section}{\Large\bfseries\color{accentblue}}{}{0em}{\thesection\quad}[\titlerule]
\titleformat{\subsection}{\large\bfseries\color{accentblue}}{}{0em}{\thesubsection\quad}
\titleformat{\subsubsection}{\normalsize\bfseries}{}{0em}{\thesubsubsection\quad}

% ============================================================
% Dokument
% ============================================================
\begin{document}

% ============================================================
% Titelseite
% ============================================================
\begin{titlepage}
  \centering
  \vspace*{2cm}
  {\Huge\bfseries\color{accentblue} KI-Copilot für Verzahnungswissen\par}
  \vspace{0.5cm}
  {\Large\color{accentgray} Vollständige Systemdokumentation\par}
  \vspace{0.3cm}
  {\large Aufbau, Datenfluss und Implementierung\par}
  \vspace{2cm}

  \begin{tcolorbox}[colback=boxbg, colframe=boxframe, width=0.85\textwidth]
    \centering\small
    Dieses Dokument beschreibt jeden Code-Baustein des modularen RAG-Systems
    (Retrieval-Augmented Generation) für technisches Verzahnungswissen --
    vom PDF-Upload bis zur Antwortgenerierung durch ein lokales Sprachmodell.
  \end{tcolorbox}

  \vspace{2cm}
  \begin{tabular}{ll}
    \textbf{Technologie-Stack:} & Python 3.9+, FastAPI, Qdrant, Ollama \\
    \textbf{Embedding-Modell:}  & BAAI/bge-m3 (sentence-transformers) \\
    \textbf{LLM:}               & llama3.2:3b (lokal via Ollama) \\
    \textbf{Vektordatenbank:}   & Qdrant (Docker, Port 6333) \\
    \textbf{Frontend:}          & HTML + Vanilla JS + Pico.css \\
  \end{tabular}

  \vfill
  {\small Stand: Mai 2026}
\end{titlepage}

% ============================================================
% Inhaltsverzeichnis
% ============================================================
\tableofcontents
\newpage

% ============================================================
\section{Systemübersicht}
% ============================================================

Das System ist ein modulares \textbf{RAG-System} (\textit{Retrieval-Augmented Generation}),
das Fachfragen zu Verzahnungen ausschließlich auf Basis hochgeladener PDF-Dokumente und
übermittelter Bauteildaten beantwortet. Es verwendet \textbf{kein} internes Vorwissen
des Sprachmodells.

\subsection{Drei Eingangsdatenströme}

\begin{enumerate}[leftmargin=2cm]
  \item \textbf{Wissensbasis:} PDF-Dokumente (Normen, Lehrbücher, Produktionsdaten),
        einmalig oder dynamisch hochgeladen, dauerhaft in Qdrant persistiert.
  \item \textbf{Bauteildaten (CAD):} Pro Anfrage übermittelte Geometrie-Metadaten
        eines konkreten Zahnrads (aktuell: Random-Stub, später: CAD-Datei-Auswertung).
  \item \textbf{Fragen:} Natürlichsprachliche Fragen des Nutzers, mehrere parallel möglich.
\end{enumerate}

\subsection{Verzeichnisstruktur}

\begin{lstlisting}[style=yamlstyle, caption={Verzeichnisstruktur des Projekts}]
KI-Copilot fuer Verzahnungswissen/
├── app/
│   ├── api/
│   │   └── main.py              # FastAPI-Endpunkte + App-Factory
│   ├── core/
│   │   ├── config.py            # Pydantic-Konfigurationsmodelle
│   │   ├── factory.py           # Komponentenaufbau aus Config
│   │   ├── interfaces.py        # Python Protocols (Schnittstellen)
│   │   ├── schema.py            # Metadatenschema-Datenstrukturen
│   │   ├── types.py             # Datenklassen (Chunk, Answer, ...)
│   │   └── utils.py             # Hilfsfunktionen (Hash, JSON)
│   ├── implementations/
│   │   ├── answer_generator_ollama.py   # Antwortgenerierung via Ollama
│   │   ├── cad_random_gear.py           # Zufaelliger Zahnrad-Generator
│   │   ├── chunker_recursive.py         # Regelbasierter Chunker
│   │   ├── chunker_semantic.py          # Semantischer Chunker
│   │   ├── embedder_bge_m3.py           # Embedding mit bge-m3
│   │   ├── metadata_extractor_ollama.py # LLM-Metadatenextraktion
│   │   ├── ollama_client.py             # HTTP-Client fuer Ollama
│   │   ├── pdf_loader_pymupdf.py        # PDF-Parsing mit PyMuPDF
│   │   ├── qdrant_store.py              # Vektordatenbank-Adapter
│   │   ├── retriever_two_stage.py       # Zweistufiger Retriever
│   │   └── text_split.py               # Satz-Tokenisierung
│   └── pipeline/
│       └── indexer.py                   # Indizierungs-Pipeline
├── frontend/
│   └── index.html               # Einzel-HTML-Frontend
├── prompts/
│   └── answer_system_prompt.txt # Prompt-Template fuer LLM
├── schemas/
│   └── gears.yaml               # Domänen-Metadatenschema
└── config.yaml                  # Zentrale Konfigurationsdatei
\end{lstlisting}

\newpage
% ============================================================
\section{Gesamtarchitektur und Datenfluss}
% ============================================================

\begin{figure}[H]
\centering
\begin{tikzpicture}[
  node distance=0.55cm and 0.9cm,
  % Eingabe-/Ausgabe-Knoten
  inbox/.style={rectangle, rounded corners=3pt, draw=accentblue, fill=boxbg,
                text width=3.4cm, align=center, minimum height=0.75cm,
                font=\small\bfseries\color{accentblue}},
  % Komponenten-Knoten mit Dateinamen
  sbox/.style={rectangle, rounded corners=3pt, draw=accentgray, fill=codebg,
               text width=3.4cm, align=center, minimum height=0.9cm,
               font=\small},
  arrow/.style={-{Stealth[length=2mm]}, thick, color=accentblue},
]

%% ---- SPALTE 1: Indizierungspfad ----
\node[inbox] (pdf)   {PDF-Upload};

\node[sbox, below=0.7cm of pdf] (loader)
  {\textbf{PDFLoader} \\ \texttt{\tiny pdf\_loader\_pymupdf.py}};

\node[sbox, below=0.5cm of loader] (chunker)
  {\textbf{Chunker} \\ \texttt{\tiny chunker\_recursive.py} \\ \texttt{\tiny chunker\_semantic.py} \\ \texttt{\tiny text\_split.py}};

\node[sbox, below=0.5cm of chunker] (metaex)
  {\textbf{MetadataExtractor} \\ \texttt{\tiny metadata\_extractor\_ollama.py} \\ \texttt{\tiny ollama\_client.py} \\ \texttt{\tiny gears.yaml}};

\node[sbox, below=0.5cm of metaex] (emb1)
  {\textbf{Embedder} \\ \texttt{\tiny embedder\_bge\_m3.py}};

\node[sbox, below=0.5cm of emb1] (qdrant)
  {\textbf{VectorStore} \\ \texttt{\tiny qdrant\_store.py}};

%% ---- SPALTE 2: CAD-Pfad ----
\node[inbox, right=1.0cm of pdf] (cad) {CAD-Daten};

\node[sbox, below=0.7cm of cad] (cadgen)
  {\textbf{CAD-Adapter} \\ \texttt{\tiny cad\_random\_gear.py}};

%% ---- SPALTE 3: Anfragepfad ----
\node[inbox, right=1.0cm of cad] (frage) {Frage(n)};

\node[sbox, below=0.7cm of frage] (emb2)
  {\textbf{Embedder} \\ \texttt{\tiny embedder\_bge\_m3.py}};

\node[sbox, below=0.5cm of emb2] (ret)
  {\textbf{TwoStageRetriever} \\ \texttt{\tiny retriever\_two\_stage.py} \\ \texttt{\tiny schema.py}};

\node[sbox, below=0.5cm of ret] (ansgen)
  {\textbf{AnswerGenerator} \\ \texttt{\tiny answer\_generator\_ollama.py} \\ \texttt{\tiny ollama\_client.py} \\ \texttt{\tiny answer\_system\_prompt.txt}};

\node[inbox, below=0.5cm of ansgen] (out)
  {Antwort + Quellen};

%% ---- Pfeile Indizierung ----
\draw[arrow] (pdf)     -- (loader);
\draw[arrow] (loader)  -- (chunker);
\draw[arrow] (chunker) -- (metaex);
\draw[arrow] (metaex)  -- (emb1);
\draw[arrow] (emb1)    -- (qdrant);

%% ---- Pfeil CAD ----
\draw[arrow] (cad) -- (cadgen);

%% ---- Pfeile Anfrage ----
\draw[arrow] (frage) -- (emb2);
\draw[arrow] (emb2)  -- (ret);

%% ---- Verbindungen zum Retriever / AnswerGenerator ----
\draw[arrow] (qdrant.east)   -- ++(0.45,0) |- (ret.west);
\draw[arrow] (cadgen.south)  -- ++(0,-0.35) -| (ret.north);
\draw[arrow] (ret)           -- (ansgen);
\draw[arrow] (cadgen.east)   -- ++(0.3,0) |- (ansgen.west);
\draw[arrow] (ansgen) -- (out);

\end{tikzpicture}
\caption{Architekturdiagramm mit Datei-Zuordnung -- jeder Knoten zeigt die zugehörigen Quelldateien}
\end{figure}

Das System besteht aus zwei Hauptpfaden:

\begin{itemize}
  \item \textbf{Indizierungspfad} (einmalig pro PDF): PDF $\to$ Text $\to$ Chunks
        $\to$ Metadaten $\to$ Vektoren $\to$ Qdrant
  \item \textbf{Anfragepfad} (pro Nutzerfrage): Frage $\to$ Embedding
        $\to$ Zweistufige Suche $\to$ LLM-Antwort
\end{itemize}

\subsection{Vollständige Datei-Zuordnung zu den Schaubild-Schritten}

Die folgende Tabelle ordnet \textbf{jede Datei aus Abschnitt~1.2} einem konkreten Schritt
im Architekturdiagramm zu. Dateien der Kernschicht (\texttt{core/}) sind querschnittlich --
sie unterstützen mehrere Schritte gleichzeitig.

\begin{longtable}{p{5.2cm} p{3.5cm} p{5.5cm}}
\toprule
\textbf{Datei} & \textbf{Schritt im Diagramm} & \textbf{Aufgabe} \\
\midrule
\endfirsthead
\toprule
\textbf{Datei} & \textbf{Schritt im Diagramm} & \textbf{Aufgabe} \\
\midrule
\endhead
\bottomrule
\endfoot

\multicolumn{3}{l}{\textit{\small Indizierungspfad}} \\[2pt]
\texttt{pdf\_loader\_pymupdf.py}     & PDFLoader            & PDF seitenweise in Text umwandeln; SHA-256 erzeugen \\[3pt]
\texttt{chunker\_recursive.py}       & Chunker              & Regelbasierte Aufteilung in Textsegmente \\[3pt]
\texttt{chunker\_semantic.py}        & Chunker              & Kosinusähnlichkeit erkennt inhaltliche Grenzen \\[3pt]
\texttt{text\_split.py}              & Chunker              & Satztrennung und Token-Approximation für beide Chunker \\[3pt]
\texttt{metadata\_extractor\_ollama.py} & MetadataExtractor & LLM extrahiert Metadaten pro Chunk als JSON \\[3pt]
\texttt{embedder\_bge\_m3.py}        & Embedder (Indizierung \& Anfrage) & BAAI/bge-m3 erzeugt Dense-Vektoren; wird in beiden Pfaden genutzt \\[3pt]
\texttt{qdrant\_store.py}            & VectorStore          & Upsert, Vektorsuche, Löschung, Dokumentliste in Qdrant \\[3pt]
\texttt{pipeline/indexer.py}         & Orchestrierung Indizierung & Koordiniert Schritte 1--6 (Loader $\to$ Qdrant) \\[6pt]

\multicolumn{3}{l}{\textit{\small CAD-Pfad}} \\[2pt]
\texttt{cad\_random\_gear.py}        & CAD-Adapter          & Generiert zufällige, geometrisch konsistente Zahnraddaten \\[6pt]

\multicolumn{3}{l}{\textit{\small Anfragepfad}} \\[2pt]
\texttt{retriever\_two\_stage.py}    & TwoStageRetriever    & Stage-1 Metadatenfilter + Stage-2 Vektorsuche \\[3pt]
\texttt{answer\_generator\_ollama.py} & AnswerGenerator     & Baut Prompt auf, ruft LLM auf, gibt Antwort + Quellen zurück \\[3pt]
\texttt{ollama\_client.py}           & MetadataExtractor \& AnswerGenerator & HTTP-Client für alle Ollama-Aufrufe \\[3pt]
\texttt{prompts/answer\_system\_prompt.txt} & AnswerGenerator & Prompt-Template mit Platzhaltern (DOMAIN, CHUNKS, FRAGE, FORMAT) \\[6pt]

\multicolumn{3}{l}{\textit{\small API- und Frontend-Schicht}} \\[2pt]
\texttt{api/main.py}                 & Gesamte Pipeline     & FastAPI-Endpunkte; startet Indizierung und parallele Anfragen \\[3pt]
\texttt{frontend/index.html}         & Nutzerinterface      & Einzel-HTML-Datei; bindet alle Endpunkte an, zeigt Antworten \\[6pt]

\multicolumn{3}{l}{\textit{\small Konfiguration und Schema}} \\[2pt]
\texttt{config.yaml}                 & Alle Schritte        & Zentrale Parameterdatei; steuert jede Komponente \\[3pt]
\texttt{schemas/gears.yaml}          & MetadataExtractor \& TwoStageRetriever & Domänenschema: Feldnamen, Typen, Filterbedingungen \\[6pt]

\multicolumn{3}{l}{\textit{\small Kernschicht -- querschnittlich in allen Schritten}} \\[2pt]
\texttt{core/types.py}               & Alle Schritte        & Gemeinsame Datenklassen (Chunk, Answer, ...) -- fließen durch die gesamte Pipeline \\[3pt]
\texttt{core/interfaces.py}          & Alle Schritte        & Protokoll-Definitionen; Vertrag zwischen Factory und Implementierungen \\[3pt]
\texttt{core/config.py}              & Startup              & Pydantic-Modelle; validiert \texttt{config.yaml} beim Serverstart \\[3pt]
\texttt{core/factory.py}             & Startup              & Instanziiert alle Komponenten aus der Config; Dependency Injection \\[3pt]
\texttt{core/schema.py}              & MetadataExtractor \& TwoStageRetriever & Lädt und repräsentiert das YAML-Domänenschema als Datenstruktur \\[3pt]
\texttt{core/utils.py}               & PDFLoader \& API     & SHA-256-Hashfunktion (Dokumenten-ID) + stabiles JSON-Serialisieren \\

\bottomrule
\caption{Zuordnung aller Dateien aus Abschnitt~1.2 zu den Schaubild-Schritten}
\end{longtable}

\newpage
% ============================================================
\section{Kernschicht: \texttt{app/core/}}
% ============================================================

Die Kernschicht enthält ausschließlich abstrakte Definitionen, Datenstrukturen und
Konfigurationslogik -- keinen Algorithmus-Code. Sie ist vollständig implementierungsunabhängig.

% ------------------------------------------------------------
\subsection{\texttt{types.py} -- Datenstrukturen}
% ------------------------------------------------------------

Definiert alle gemeinsamen Datenklassen des Systems als unveränderliche
\texttt{@dataclass(frozen=True)}-Objekte. Diese fließen als ``Datenträger'' durch
alle Komponenten.

\begin{lstlisting}[style=pythonstyle, caption={Kerntypen des Systems (types.py)}]
@dataclass(frozen=True)
class RawDocumentPage:
    page_number: int   # 1-basiert
    text: str          # Extrahierter Seitentext

@dataclass(frozen=True)
class RawDocument:
    source_path: str   # Absoluter Pfad zur PDF
    doc_hash: str      # SHA-256 der Datei (zur Identifikation)
    pages: list[RawDocumentPage]

@dataclass(frozen=True)
class Chunk:
    text: str          # Textsegment
    source_path: str
    page_number: int
    position: int      # Monoton steigende Position im Dokument
    doc_hash: str      # Zugehoerige Dokument-ID

@dataclass(frozen=True)
class EmbeddedChunk:
    chunk: Chunk
    dense_vector: list[float]    # Embedding-Vektor
    sparse_vector: Optional[...]  # fuer spaeteren Hybrid-Modus
    metadata: dict[str, Any]     # LLM-extrahierte Metadaten

@dataclass(frozen=True)
class RetrievedChunk:
    chunk: Chunk
    metadata: dict[str, Any]
    similarity: float   # Aehnlichkeitsscore aus Stage 2
\end{lstlisting}

\begin{infobox}[Designentscheidung: frozen=True]
Alle Datenklassen sind eingefroren (\texttt{frozen=True}). Das verhindert versehentliche
Mutationen, macht die Objekte hashbar und thread-sicher -- wichtig, da der Server
parallele Anfragen über \texttt{asyncio.gather} verarbeitet.
\end{infobox}

Das \texttt{Answer}-Ergebnis ist als \texttt{TypedDict} definiert, da es direkt
als JSON-serialisierbare Struktur an FastAPI zurückgegeben wird:

\begin{lstlisting}[style=pythonstyle, caption={Antwortformat (types.py)}]
class AnswerSource(TypedDict):
    qid: str           # "Q1", "Q2", ... (Quellenkennung im Antworttext)
    source_path: str
    page_number: int
    similarity: float
    text: str          # Chunk-Originaltext

class Answer(TypedDict):
    question: str
    answer_text: str
    sources: list[AnswerSource]
\end{lstlisting}

% ------------------------------------------------------------
\subsection{\texttt{interfaces.py} -- Komponentenschnittstellen}
% ------------------------------------------------------------

Definiert alle Komponenten als Python \texttt{Protocol}-Klassen. Protocols ermöglichen
strukturelles Subtyping: Eine Klasse implementiert eine Schnittstelle, wenn sie die
passenden Methoden hat -- ohne explizite Vererbung.

\begin{lstlisting}[style=pythonstyle, caption={Protokoll-Schnittstellen (interfaces.py)}]
class DocumentLoader(Protocol):
    def load(self, file_path: Path) -> RawDocument: ...

class Chunker(Protocol):
    def chunk(self, document: RawDocument) -> list[Chunk]: ...

class Embedder(Protocol):
    def embed(self, texts: list[str]) -> EmbeddingResult: ...

class MetadataExtractor(Protocol):
    def extract(self, chunk: Chunk, schema: MetadataSchema) -> dict: ...

class CADAdapter(Protocol):
    def extract(self, file_path: Optional[Path]) -> dict: ...

class VectorStore(Protocol):
    def upsert(self, chunks: list[EmbeddedChunk]) -> None: ...
    def search(self, query_vector, *, filter, top_k, threshold) -> list[SearchResult]: ...
    def delete_by_doc_hash(self, doc_hash: str) -> None: ...
    def list_documents(self) -> list[DocumentInfo]: ...

class Retriever(Protocol):
    def retrieve(self, question: str, cad_metadata: dict) -> list[RetrievedChunk]: ...

class AnswerGenerator(Protocol):
    def generate(self, *, question, chunks, cad_metadata, output_format) -> Answer: ...
\end{lstlisting}

\begin{infobox}[Warum Protocols statt ABCs?]
\texttt{Protocol} (PEP 544) erlaubt es, Implementierungen auszutauschen ohne gemeinsame
Basisklasse. Ein zukünftiger \texttt{OpenAIEmbedder} muss nur die Methode \texttt{embed}
haben -- er muss nirgendwo von \texttt{Embedder} erben. Das senkt den Kopplungsgrad.
\end{infobox}

% ------------------------------------------------------------
\subsection{\texttt{config.py} -- Konfigurationsmodelle}
% ------------------------------------------------------------

Alle Konfigurationsparameter sind als Pydantic-\texttt{BaseModel}-Klassen typisiert und
werden beim Serverstart gegen \texttt{config.yaml} validiert.

\begin{lstlisting}[style=pythonstyle, caption={Konfigurationsstruktur (config.py)}]
class AppConfig(BaseModel):
    domain:            DomainConfig
    embedder:          EmbedderConfig
    chunker:           ChunkerConfig
    metadata_extractor: MetadataExtractorConfig
    cad_adapter:       CADAdapterConfig
    vector_store:      VectorStoreConfig
    retriever:         RetrieverConfig
    answer_generator:  AnswerGeneratorConfig
    frontend:          FrontendConfig

def load_config(path: Path) -> AppConfig:
    raw = yaml.safe_load(path.read_text(encoding="utf-8"))
    return AppConfig.model_validate(raw)
\end{lstlisting}

Jeder Abschnitt der \texttt{config.yaml} wird einem eigenen Modell zugeordnet. Pydantic
wirft beim Start einen Validierungsfehler, wenn ein Pflichtfeld fehlt oder ein Wert
außerhalb des erlaubten Wertebereichs liegt (z.\,B. \texttt{device} $\notin$ \{cuda, cpu, mps\}).

\vspace{0.5em}
\begin{table}[H]
\centering
\small
\begin{tabularx}{\textwidth}{lX}
\toprule
\textbf{Config-Sektion} & \textbf{Wichtigste Parameter} \\
\midrule
\texttt{domain} & Name, Pfad zu Schema-YAML, Pfad zu Prompt-Datei \\
\texttt{embedder} & Modellname, Gerät (mps/cuda/cpu), max. Sequenzlänge \\
\texttt{chunker} & Implementierung (semantic/recursive), Schwellenwert, Min/Max-Tokens \\
\texttt{metadata\_extractor} & Ollama-Modell, URL, Timeout, max. Wiederholungsversuche \\
\texttt{retriever} & stage1\_strict, stage1\_relax\_on\_empty, top\_k, min\_similarity \\
\texttt{answer\_generator} & Modell, Standardformat, max. Tokens, Temperatur \\
\bottomrule
\end{tabularx}
\caption{Konfigurationssektionen und ihre Parameter}
\end{table}

% ------------------------------------------------------------
\subsection{\texttt{schema.py} -- Metadatenschema}
% ------------------------------------------------------------

Definiert die Datenstruktur für das domänenspezifische Metadatenschema, das aus einer
externen YAML-Datei geladen wird.

\begin{lstlisting}[style=pythonstyle, caption={Schema-Datenstrukturen (schema.py)}]
@dataclass(frozen=True)
class SchemaField:
    name: str
    type: str                      # "string", "number", ...
    filter_type: Optional[FilterType]  # "exact", "range", "set"
    enum: Optional[list[str]]      # erlaubte Werte
    nullable: bool
    description: Optional[str]
    range_fields: Optional[list[str]]  # z.B. ["modul_min", "modul_max"]

@dataclass(frozen=True)
class MetadataSchema:
    domain: str
    description: Optional[str]
    fields: list[SchemaField]
    extraction_prompt_hint: Optional[str]

    @property
    def filter_fields(self) -> list[SchemaField]:
        return [f for f in self.fields if f.filter_type is not None]
\end{lstlisting}

% ------------------------------------------------------------
\subsection{\texttt{utils.py} -- Hilfsfunktionen}
% ------------------------------------------------------------

\begin{lstlisting}[style=pythonstyle, caption={Hilfsfunktionen (utils.py)}]
def sha256_file(path: Path) -> str:
    # SHA-256-Hash einer Datei (chunked, speichereffizient)
    h = hashlib.sha256()
    with path.open("rb") as f:
        for chunk in iter(lambda: f.read(1024 * 1024), b""):
            h.update(chunk)
    return h.hexdigest()

def stable_json_dumps(obj: Any) -> str:
    # Deterministisches JSON (sort_keys) fuer Logs
    return json.dumps(obj, ensure_ascii=False, sort_keys=True, indent=2)
\end{lstlisting}

Der SHA-256-Hash dient als stabiler Dokumentenidentifikator (\texttt{doc\_hash}).
Wird dieselbe Datei erneut hochgeladen, ist der Hash identisch -- das ermöglicht
duplicate detection und zuverlässige Löschung aller zugehörigen Chunks.

% ------------------------------------------------------------
\subsection{\texttt{factory.py} -- Komponentenaufbau}
% ------------------------------------------------------------

Die Factory liest die validierte Konfiguration und instanziiert alle konkreten
Implementierungen. Die gesamte Abhängigkeitsauflösung findet hier statt.

\begin{lstlisting}[style=pythonstyle, caption={Factory-Funktion (factory.py)}]
def build_components(config: AppConfig, *, base_dir: Path) -> Components:
    loader   = PDFLoader()
    embedder = BGEM3Embedder(model_name=..., device=..., ...)

    # Auswahl Chunker-Implementierung
    if config.chunker.implementation == "semantic":
        chunker = SemanticChunker(embedder=embedder, threshold=..., ...)
    elif config.chunker.implementation == "recursive":
        chunker = RecursiveTextChunker(...)

    metadata_extractor = OllamaMetadataExtractor(...)
    cad_adapter        = RandomGearGenerator()
    vector_store       = QdrantStore(host=..., port=..., ...)
    retriever          = TwoStageRetriever(embedder=embedder, store=vector_store, ...)
    answer_generator   = OllamaAnswerGenerator(...)

    return Components(loader, chunker, embedder, ...)
\end{lstlisting}

\begin{infobox}[Austauschbarkeit durch Factory]
Will man z.\,B. den Chunker von \texttt{recursive} auf \texttt{semantic} umstellen,
genügt eine Änderung in \texttt{config.yaml}. Die Factory instanziiert dann automatisch
die andere Klasse. Kein Code-Eingriff notwendig.
\end{infobox}

\newpage
% ============================================================
\section{Indizierungspipeline: \texttt{app/pipeline/indexer.py}}
% ============================================================

Der \texttt{KnowledgeBaseIndexer} orchestriert den vollständigen Weg von der PDF-Datei
bis zu den in Qdrant gespeicherten Vektoren.

\begin{lstlisting}[style=pythonstyle, caption={Indizierungs-Pipeline (indexer.py)}]
def index_pdf(self, file_path: Path) -> DocumentInfo:
    # 1. PDF laden -> RawDocument
    doc = self.loader.load(file_path)

    # 2. In Chunks aufteilen -> list[Chunk]
    chunks = self.chunker.chunk(doc)

    # 3. Pro Chunk: Metadaten per LLM extrahieren -> list[dict]
    meta_list = []
    for c in chunks:
        raw_meta = self.metadata_extractor.extract(c, self.schema)
        meta_list.append(_sanitize_metadata(self.schema, raw_meta))

    # 4. Alle Chunk-Texte gemeinsam einbetten -> list[list[float]]
    vectors = self.embedder.embed([c.text for c in chunks]).dense_vectors

    # 5. EmbeddedChunks zusammenstellen
    embedded = [EmbeddedChunk(chunk=c, dense_vector=v, sparse_vector=None,
                              metadata=m) for c, v, m in zip(chunks, vectors, meta_list)]

    # 6. In Qdrant speichern
    self.store.upsert(embedded)
    return DocumentInfo(source_path=..., doc_hash=..., chunk_count=len(embedded))
\end{lstlisting}

\begin{table}[H]
\centering
\small
\begin{tabularx}{\textwidth}{clX}
\toprule
\textbf{Schritt} & \textbf{Komponente} & \textbf{Ausgabe} \\
\midrule
1 & PDFLoader & \texttt{RawDocument} (Seiten mit Text + Hash) \\
2 & Chunker & \texttt{list[Chunk]} (Textsegmente mit Positionsinfo) \\
3 & MetadataExtractor & \texttt{list[dict]} (LLM-extrahierte Metadaten) \\
4 & Embedder & \texttt{list[list[float]]} (Dichte Vektoren) \\
5 & --- & \texttt{list[EmbeddedChunk]} (Chunks + Vektoren + Metadaten) \\
6 & VectorStore & gespeichert in Qdrant \\
\bottomrule
\end{tabularx}
\caption{Schritte der Indizierungspipeline}
\end{table}

Die Funktion \texttt{\_sanitize\_metadata} filtert LLM-Ausgaben: Nur Felder, die im
Schema definiert sind, werden in Qdrant gespeichert. Das verhindert das Einschleusen
unbekannter Schlüssel.

\newpage
% ============================================================
\section{Implementierungen: \texttt{app/implementations/}}
% ============================================================

% ------------------------------------------------------------
\subsection{PDF-Loader: \texttt{pdf\_loader\_pymupdf.py}}
% ------------------------------------------------------------

\begin{lstlisting}[style=pythonstyle, caption={PDF-Loader mit PyMuPDF}]
class PDFLoader:
    def load(self, file_path: Path) -> RawDocument:
        doc_hash = sha256_file(file_path)  # Datei-Fingerabdruck
        doc = fitz.open(file_path)         # fitz = PyMuPDF
        pages = []
        for i in range(doc.page_count):
            page = doc.load_page(i)
            text = page.get_text("text") or ""
            pages.append(RawDocumentPage(page_number=i + 1, text=text))
        doc.close()
        return RawDocument(source_path=str(file_path),
                           doc_hash=doc_hash, pages=pages)
\end{lstlisting}

PyMuPDF (\texttt{fitz}) extrahiert den Reintext seitenweise. Die Seitennummer wird
1-basiert gespeichert und bleibt bis zur Antwort erhalten -- so kann jede Quellenangabe
die genaue Seitenzahl nennen.

% ------------------------------------------------------------
\subsection{Text-Splitting: \texttt{text\_split.py}}
% ------------------------------------------------------------

Hilfsfunktionen für beide Chunker-Implementierungen:

\begin{lstlisting}[style=pythonstyle, caption={Satztrennung und Token-Approximation}]
# Satzerkennung: Trennt nach [.!?] + Leerzeichen vor Grossbuchstabe/Zahl
_sentence_re = re.compile(r"(?<=[.!?])\s+(?=[A-ZAEOEU0-9])")

def split_sentences(text: str) -> list[str]:
    parts = _sentence_re.split(text)
    return [p.strip() for p in parts if p.strip()]

def approx_tokens(text: str) -> int:
    # Wortanzahl als Naherung (modell-/bibliotheksunabhaengig)
    return max(1, len(re.findall(r"\S+", text)))
\end{lstlisting}

Die Token-Approximation über Wortzählung ist bewusst modellunabhängig -- die
Konfigurationsparameter \texttt{min\_chunk\_size} und \texttt{max\_chunk\_size}
beziehen sich auf diese Einheit.

% ------------------------------------------------------------
\subsection{Semantischer Chunker: \texttt{chunker\_semantic.py}}
% ------------------------------------------------------------

Der semantische Chunker erkennt inhaltliche Grenzen zwischen Sätzen über den
Kosinus-Abstand ihrer Embeddings.

\begin{lstlisting}[style=pythonstyle, caption={Semantische Grenzenerkennung}]
def chunk(self, document: RawDocument) -> list[Chunk]:
    for page in document.pages:
        sentences = split_sentences(page.text)
        sent_vecs = self.embedder.embed(sentences).dense_vectors

        # Grenzen dort, wo Kosinusaehnlichkeit unter Schwellenwert faellt
        boundaries = set()
        for i in range(1, len(sentences)):
            sim = _cosine(sent_vecs[i-1], sent_vecs[i])
            if sim < self.threshold:  # Default: 0.75
                boundaries.add(i)

        # Segmente zu Chunks packen (Max-Token-Limit beachten, Overlap hinzufuegen)
        ...
\end{lstlisting}

\begin{table}[H]
\centering
\small
\begin{tabular}{ll}
\toprule
\textbf{Parameter} & \textbf{Bedeutung} \\
\midrule
\texttt{threshold = 0.75} & Mindest-Ähnlichkeit, unter der eine Grenze gesetzt wird \\
\texttt{min\_chunk\_tokens = 100} & Chunks unter dieser Größe werden verworfen \\
\texttt{max\_chunk\_tokens = 512} & Maximale Chunk-Größe (Tokens) \\
\texttt{overlap\_sentences = 1} & Sätze, die in aufeinanderfolgenden Chunks überlappen \\
\bottomrule
\end{tabular}
\caption{Parameter des semantischen Chunkers}
\end{table}

% ------------------------------------------------------------
\subsection{Rekursiver Chunker: \texttt{chunker\_recursive.py}}
% ------------------------------------------------------------

Einfachere Alternative ohne LLM-Aufruf beim Chunking: Sätze werden sequentiell
zusammengepackt bis das Token-Limit erreicht ist.

\begin{lstlisting}[style=pythonstyle, caption={Regelbasierter Chunker}]
def chunk(self, document: RawDocument) -> list[Chunk]:
    for page in document.pages:
        sentences = split_sentences(page.text)
        start = 0
        while start < len(sentences):
            cur, cur_tokens = [], 0
            i = start
            while i < len(sentences):
                s_tokens = approx_tokens(sentences[i])
                if cur and (cur_tokens + s_tokens) > self.max_chunk_tokens:
                    break  # naechsten Chunk beginnen
                cur.append(sentences[i])
                cur_tokens += s_tokens
                i += 1
            # Overlap: naechster Chunk startet `overlap_sentences` Saetze zurueck
            start = max(i - self.overlap_sentences, start + 1)
\end{lstlisting}

\begin{infobox}[Wahl des Chunkers]
Der rekursive Chunker ist schneller (kein Embedding-Aufruf), produziert aber gleichmäßigere,
inhaltlich weniger kohärente Chunks. Der semantische Chunker ist langsamer, erzeugt
jedoch thematisch konsistentere Segmente -- empfohlen für Fachliteratur.
In \texttt{config.yaml}: \texttt{chunker.implementation: "semantic"} oder \texttt{"recursive"}.
\end{infobox}

% ------------------------------------------------------------
\subsection{Embedder: \texttt{embedder\_bge\_m3.py}}
% ------------------------------------------------------------

\begin{lstlisting}[style=pythonstyle, caption={BGE-M3 Embedding-Wrapper}]
class BGEM3Embedder:
    def __init__(self, *, model_name, device, max_length, use_sparse):
        self._model = SentenceTransformer(model_name, device=device)

    def embed(self, texts: list[str]) -> EmbeddingResult:
        vectors = self._model.encode(
            texts,
            normalize_embeddings=True,   # L2-Normalisierung fuer Kosinus-Aehnlichkeit
            convert_to_numpy=False,
            show_progress_bar=False,
        )
        dense_vectors = [list(map(float, v)) for v in vectors]
        return EmbeddingResult(dense_vectors=dense_vectors, sparse_vectors=None)
\end{lstlisting}

Das Modell \texttt{BAAI/bge-m3} ist mehrsprachig und unterstützt bis zu 8192 Tokens
pro Sequenz. Die L2-Normalisierung (\texttt{normalize\_embeddings=True}) stellt sicher,
dass der Kosinus-Abstand korrekt als Ähnlichkeitsmaß verwendet werden kann.

\begin{definitionbox}[Wichtig: Gleiches Modell überall]
Wissensbasis und Anfrage müssen denselben Embedder verwenden, da sie im gleichen
Vektorraum repräsentiert sein müssen. Die Factory stellt sicher, dass beide
Codepfade dieselbe \texttt{BGEM3Embedder}-Instanz verwenden.
\end{definitionbox}

% ------------------------------------------------------------
\subsection{Metadaten-Extraktor: \texttt{metadata\_extractor\_ollama.py}}
% ------------------------------------------------------------

Für jeden Chunk wird ein LLM (llama3.2:3b) aufgerufen, um domänenspezifische
Metadaten im JSON-Format zu extrahieren.

\begin{lstlisting}[style=pythonstyle, caption={LLM-basierte Metadatenextraktion}]
def extract(self, chunk: Chunk, schema: MetadataSchema) -> dict:
    # Prompt mit Schema-Beschreibung und Chunk-Text aufbauen
    prompt = (
        "Gib AUSSCHLIESSLICH ein JSON-Objekt zurueck.\n"
        f"SCHEMA:\n{schema_desc}\n\n"
        f"CHUNK:\n{chunk.text}"
    )
    # Bis zu max_retries Versuche (LLM-Ausgabe kann ungueltig sein)
    for _ in range(self.max_retries):
        try:
            data = self.client.generate_json(model=..., prompt=...,
                                             temperature=0.0, max_tokens=400)
            if isinstance(data, dict):
                return data
        except Exception:
            continue
    return {}  # Fallback: leere Metadaten
\end{lstlisting}

\begin{table}[H]
\centering
\small
\begin{tabular}{lp{9cm}}
\toprule
\textbf{Metadatenfeld} & \textbf{Beschreibung (aus gears.yaml)} \\
\midrule
\texttt{verzahnungstyp} & Art der Verzahnung: Stirnrad, Schrägverzahnung, Kegelrad, ... \\
\texttt{modul} & Modul in mm (Bereichsfeld: modul\_min, modul\_max) \\
\texttt{norm} & Referenznorm (DIN 3960, ISO 1328 o.\,ä.) \\
\texttt{themenkategorie} & Geometrie, Festigkeit, Fertigung, Werkstoff, ... \\
\texttt{anwendungsbereich} & Getriebe, Antrieb, Lenkung, ... (optional) \\
\bottomrule
\end{tabular}
\caption{Metadatenfelder für die Domäne Verzahnungen (schemas/gears.yaml)}
\end{table}

% ------------------------------------------------------------
\subsection{CAD-Adapter: \texttt{cad\_random\_gear.py}}
% ------------------------------------------------------------

Aktuell generiert dieser Adapter zufällige, geometrisch konsistente Zahnraddaten.
Er implementiert dieselbe \texttt{CADAdapter}-Schnittstelle wie ein späterer echter
CAD-File-Parser.

\begin{lstlisting}[style=pythonstyle, caption={Zufalls-Zahnrad-Generator}]
class RandomGearGenerator:
    def extract(self, file_path=None) -> dict:
        typ   = random.choice(["Stirnrad", "Schraegverzahnung"])
        modul = random.choice([1.0, 1.5, 2.0, 2.5, 3.0, 4.0, 5.0, 6.0])
        z     = random.randint(15, 60)
        alpha = 20.0           # Standardeingriffswinkel
        beta  = 0.0 if typ == "Stirnrad" else random.choice([8,10,12,15,20])
        x     = round(random.uniform(-0.3, 0.5), 2)  # Profilverschiebung

        d  = modul * z                        # Teilkreisdurchmesser
        da = d + 2 * modul * (1 + x)          # Kopfkreisdurchmesser
        df = d - 2 * modul * (1.25 - x)       # Fusskreisdurchmesser

        return {"verzahnungstyp": typ, "modul": modul, "zaehnezahl": z,
                "eingriffswinkel": alpha, "schraegungswinkel": beta,
                "profilverschiebung": x, "teilkreisdurchmesser": round(d,2),
                "kopfkreisdurchmesser": round(da,2), "fusskreisdurchmesser": round(df,2),
                "zahnbreite": round(modul * random.uniform(8, 14), 1),
                "werkstoff": random.choice(["16MnCr5","20MnCr5","42CrMo4","C45"]),
                "haerte": random.choice(["verguetet","einsatzgehaertet","nitriert"]),
                "verzahnungsqualitaet": random.choice([6, 7, 8])}
\end{lstlisting}

Die geometrischen Zusammenhänge (Kopfkreis, Fußkreis) sind nach DIN 3960 berechnet
und in sich konsistent -- die generierten Daten sind technisch plausibel.

% ------------------------------------------------------------
\subsection{Qdrant-VectorStore: \texttt{qdrant\_store.py}}
% ------------------------------------------------------------

Adapter zwischen dem System und der Qdrant-Vektordatenbank. Verwaltet Collection,
Upsert, Suche und Löschung.

\begin{lstlisting}[style=pythonstyle, caption={Vektordatenbank-Adapter (Qdrant)}]
class QdrantStore:
    def __init__(self, *, host, port, collection_name):
        self.client = QdrantClient(host=host, port=port)

    def upsert(self, chunks: list[EmbeddedChunk]) -> None:
        self._ensure_collection(vector_size=len(chunks[0].dense_vector))
        points = []
        for ec in chunks:
            # Stabiler Punkt-ID aus (doc_hash, position) -- kein UUID-Generieren
            pid = abs(hash((ec.chunk.doc_hash, ec.chunk.position))) % (2**63)
            payload = {"text": ..., "source_path": ..., "page_number": ...,
                       "position": ..., "doc_hash": ..., "metadata": ec.metadata}
            points.append(PointStruct(id=pid, vector=ec.dense_vector, payload=payload))
        self.client.upsert(collection_name=..., points=points)

    def search(self, query_vector, *, filter, top_k, threshold):
        qfilter = _dict_filter_to_qdrant(filter)
        response = self.client.query_points(   # Qdrant >= 1.10
            collection_name=self.collection_name,
            query=query_vector,
            query_filter=qfilter,
            limit=top_k,
            score_threshold=threshold,
            with_payload=True,
        )
        return [SearchResult(chunk=..., score=h.score, ...) for h in response.points]
\end{lstlisting}

\begin{infobox}[Stabiler Punkt-ID]
Statt einer zufälligen UUID wird der Punkt-ID als deterministischer Hash aus
\texttt{(doc\_hash, position)} berechnet. Bei erneutem Hochladen derselben Datei
werden die alten Punkte überschrieben (\texttt{upsert}), nicht dupliziert.
\end{infobox}

Die Filterfunktion \texttt{\_dict\_filter\_to\_qdrant} übersetzt das interne
Dict-Format in Qdrant-\texttt{Filter}-Objekte und unterstützt drei Konditionstypen:

\begin{table}[H]
\centering
\small
\begin{tabular}{llp{7cm}}
\toprule
\textbf{Typ} & \textbf{Qdrant-Kondition} & \textbf{Verwendung} \\
\midrule
\texttt{match} & \texttt{FieldCondition(match=MatchValue)} & Exakter Wertvergleich (z.\,B. Verzahnungstyp) \\
\texttt{range} & \texttt{FieldCondition(range=Range)} & Wertebereich (z.\,B. Modul) \\
\texttt{contains} & \texttt{FieldCondition(match=MatchValue)} & Listenzugehörigkeit (z.\,B. Themenkategorie) \\
\bottomrule
\end{tabular}
\caption{Unterstützte Filterkonditionstypen}
\end{table}

Alle Konditionstypen unterstützen \texttt{or\_empty=True}: Ein Chunk ohne dieses
Metadatenfeld wird nicht ausgeschlossen.

% ------------------------------------------------------------
\subsection{Zweistufiger Retriever: \texttt{retriever\_two\_stage.py}}
% ------------------------------------------------------------

Das Herzstück des RAG-Systems. Findet für eine Frage die relevantesten Chunks
durch einen deterministischen Metadatenfilter (Stage 1) gefolgt von einer
Vektorähnlichkeitssuche (Stage 2).

\subsubsection{Stage 1: Deterministischer Metadatenfilter}

Baut einen Qdrant-Filter aus den CAD-Metadaten des Bauteils:

\begin{lstlisting}[style=pythonstyle, caption={Stage-1-Filter-Aufbau}]
def _build_stage1_filter(schema, cad, *, relax_factor=1.0) -> dict:
    must = []
    for field in schema.filter_fields:
        if field.filter_type == "exact":
            # z.B.: Verzahnungstyp des Bauteils muss dem Chunk-Typ entsprechen
            must.append({"key": f"metadata.{field.name}",
                         "match": cad.get(field.name), "or_empty": True})
        elif field.filter_type == "range":
            # z.B.: Bauteil-Modul muss im Modul-Bereich des Chunks liegen
            # relax_factor erweitert den Bereich bei zu wenig Treffern
            hi = float(cad_value) * relax_factor
            lo = float(cad_value) / relax_factor
            must.append({"key": f"metadata.{min_key}", "range": {"lte": hi}, "or_empty": True})
            must.append({"key": f"metadata.{max_key}", "range": {"gte": lo}, "or_empty": True})
    return {"must": must}
\end{lstlisting}

\subsubsection{Stage 2: Embedding-Vektorsuche + Relax-Logik}

\begin{lstlisting}[style=pythonstyle, caption={Zweistufiger Retrieval-Ablauf}]
def retrieve(self, question: str, cad_metadata: dict) -> list[RetrievedChunk]:
    # Frage einbetten
    qvec = self.embedder.embed([question]).dense_vectors[0]

    # Stage 1: Qdrant-Suche mit Metadatenfilter
    filter0 = _build_stage1_filter(self.schema, cad_metadata, relax_factor=1.0)
    hits = self.store.search(qvec, filter=filter0,
                             top_k=self.top_k, threshold=self.min_similarity)

    # Relax-Logik: zu wenige Treffer -> Filter lockern
    if self.stage1_relax_on_empty and len(hits) < self.stage1_min_candidates:
        for factor in (1.5, 2.0):
            hits = self.store.search(qvec,
                filter=_build_stage1_filter(self.schema, cad_metadata, relax_factor=factor),
                top_k=self.top_k, threshold=self.min_similarity)
            if hits:
                break

    return [RetrievedChunk(chunk=h.chunk, metadata=h.metadata, similarity=h.score)
            for h in hits]
\end{lstlisting}

\begin{figure}[H]
\centering
\begin{tikzpicture}[
  node distance=0.5cm and 1.5cm,
  box/.style={rectangle, rounded corners, draw=accentblue, fill=boxbg,
              text width=5.5cm, align=center, minimum height=0.9cm, font=\small},
  dbox/.style={rectangle, rounded corners, draw=accentgray, fill=codebg,
               text width=5.5cm, align=center, minimum height=0.9cm, font=\small},
  arrow/.style={-{Stealth[length=2mm]}, thick, color=accentblue},
]
\node[box]  (q)     {1. Frage einbetten (bge-m3)};
\node[box,  below=0.4cm of q]  (s1)    {2. Stage-1-Filter aufbauen (relax=1.0)};
\node[box,  below=0.4cm of s1] (qdrant){3. Qdrant-Suche mit Filter};
\node[dbox, below=0.4cm of qdrant] (check) {Treffer $\geq$ min\_candidates?};
\node[box,  below=0.4cm of check]  (relax) {4a. Falls Nein: Filter lockern (1.5x, 2.0x)};
\node[box,  below=0.4cm of relax]  (topk)  {4b. Top-k Chunks (Score $\geq$ 0.65)};
\node[box,  below=0.4cm of topk]   (out)   {5. list[RetrievedChunk] ausgeben};

\draw[arrow] (q)      -- (s1);
\draw[arrow] (s1)     -- (qdrant);
\draw[arrow] (qdrant) -- (check);
\draw[arrow] (check)  -- (relax);
\draw[arrow] (relax)  -- (topk);
\draw[arrow] (topk)   -- (out);
\end{tikzpicture}
\caption{Ablauf des zweistufigen Retrievers}
\end{figure}

% ------------------------------------------------------------
\subsection{Antwortgenerator: \texttt{answer\_generator\_ollama.py}}
% ------------------------------------------------------------

\begin{lstlisting}[style=pythonstyle, caption={Prompt-Aufbau und LLM-Aufruf}]
def generate(self, *, question, chunks, cad_metadata, output_format) -> Answer:
    # Chunk-Texte als nummerierte Quellenblock aufbauen
    chunk_lines = []
    sources = []
    for idx, rc in enumerate(chunks, start=1):
        qid = f"Q{idx}"
        chunk_lines.append(f"[{qid}] Quelle: {Path(rc.chunk.source_path).name},
                             Seite {rc.chunk.page_number}")
        chunk_lines.append(rc.chunk.text.strip())
        chunk_lines.append("---")
        sources.append({"qid": qid, "source_path": ..., "similarity": ...})

    # Prompt-Template aus Datei befuellen
    prompt = self.prompt_template.format(
        DOMAIN=self.domain_name,
        CAD_METADATA_JSON=stable_json_dumps(cad_metadata),
        CHUNKS_BLOCK="\n".join(chunk_lines),
        QUESTION=question,
        FORMAT_INSTRUCTION=_FORMAT_INSTRUCTIONS[fmt],
    )

    answer_text = self.client.generate(model=self.model_name, prompt=prompt, ...)
    return {"question": question, "answer_text": answer_text, "sources": sources}
\end{lstlisting}

\subsubsection{Prompt-Template (\texttt{prompts/answer\_system\_prompt.txt})}

Das Prompt-Template erzwingt striktes RAG -- das LLM darf nur die übergebenen
Chunks und Bauteildaten als Quellen verwenden:

\begin{lstlisting}[style=yamlstyle, caption={System-Prompt-Struktur}]
Du bist ein technischer Assistent fuer die Domaene {DOMAIN}.

REGELN:
1. Antworte AUSSCHLIESSLICH auf Basis der Wissensauszuege und Bauteildaten.
2. Verwende KEIN externes Vorwissen.
3. Wenn Information fehlt: Sage das explizit.
4. Haenge JEDE inhaltliche Aussage eine Quellenmarkierung [Q1], [Q2]... an.

BAUTEILDATEN: {CAD_METADATA_JSON}
WISSENSAUSZUEGE: {CHUNKS_BLOCK}
FRAGE: {QUESTION}
AUSGABEFORMAT: {FORMAT_INSTRUCTION}

ANTWORT:
\end{lstlisting}

\begin{table}[H]
\centering
\small
\begin{tabular}{lp{9cm}}
\toprule
\textbf{Format} & \textbf{Anweisung an das LLM} \\
\midrule
\texttt{kurz} & Fließtext, prägnant, max. 100 Wörter \\
\texttt{standard} & Strukturiert mit Zwischenüberschriften, ca. 200 Wörter \\
\texttt{ausführlich} & Mit Zwischenüberschriften, ca. 400 Wörter \\
\texttt{stichpunkte} & Stichpunktliste, max. 7 Punkte \\
\texttt{tabellarisch} & Tabellarisch, wenn Vergleichsdaten gefragt sind \\
\bottomrule
\end{tabular}
\caption{Ausgabeformate und ihre Prompt-Anweisungen}
\end{table}

% ------------------------------------------------------------
\subsection{Ollama-HTTP-Client: \texttt{ollama\_client.py}}
% ------------------------------------------------------------

\begin{lstlisting}[style=pythonstyle, caption={HTTP-Client fuer Ollama}]
class OllamaClient:
    def generate(self, *, model, prompt, temperature=None, max_tokens=None) -> str:
        payload = {"model": model, "prompt": prompt, "stream": False}
        if temperature is not None: payload["temperature"] = temperature
        if max_tokens is not None:  payload["num_predict"] = max_tokens

        with httpx.Client(timeout=self.timeout_s) as client:
            r = client.post(f"{self.base_url}/api/generate", json=payload)
            r.raise_for_status()
            return str(r.json().get("response", "")).strip()

    def generate_json(self, ...) -> Any:
        text = self.generate(...)
        try:
            return json.loads(text)
        except Exception:
            # Fallback: JSON-Objekt aus umgebendem Text extrahieren
            start, end = text.find("{"), text.rfind("}")
            if start != -1 and end > start:
                return json.loads(text[start:end+1])
            raise
\end{lstlisting}

\texttt{generate\_json} wird für die Metadatenextraktion verwendet: Es versucht,
ein JSON-Objekt aus der LLM-Ausgabe zu parsen, auch wenn das Modell die Antwort
in Prosa einbettet.

\newpage
% ============================================================
\section{API-Schicht: \texttt{app/api/main.py}}
% ============================================================

FastAPI-Server mit vier Endpunkten. Die App wird durch die Factory-Funktion
\texttt{create\_app()} aufgebaut.

\begin{lstlisting}[style=pythonstyle, caption={API-Endpunkte (Uebersicht)}]
@app.get("/")             # Liefert das Frontend (index.html)
@app.post("/upload")      # PDF hochladen und indizieren
@app.get("/documents")    # Liste aller indizierten Dokumente
@app.delete("/documents/{doc_hash}")  # Dokument entfernen
@app.get("/cad/random")   # Zufaelliges Zahnrad generieren
@app.post("/ask")         # Fragen stellen und Antworten erhalten
\end{lstlisting}

\subsection{Request/Response-Modelle}

\begin{lstlisting}[style=pythonstyle, caption={Pydantic-Modelle fuer /ask}]
class AskRequest(BaseModel):
    questions: list[str] = Field(min_length=1)  # mind. 1 Frage
    cad_metadata: dict[str, Any] = Field(default_factory=dict)
    format: Optional[OutputFormat] = None       # kurz/standard/...

class AskResponse(BaseModel):
    cad_metadata: dict[str, Any]
    answers: list[Answer]
\end{lstlisting}

\subsection{Parallele Frageverarbeitung}

\begin{lstlisting}[style=pythonstyle, caption={Parallele Anfrageverarbeitung mit asyncio}]
@app.post("/ask", response_model=AskResponse)
async def ask(req: AskRequest):
    questions = [q.strip() for q in req.questions if q.strip()]
    fmt = req.format or config.answer_generator.default_format
    cad = req.cad_metadata or {}

    # Alle Fragen PARALLEL verarbeiten (jede eigene Pipeline)
    tasks = [_answer_one(q, cad, fmt) for q in questions]
    try:
        answers = await asyncio.gather(*tasks)
    except Exception as e:
        logger.exception("ask_failed")
        raise HTTPException(status_code=500, detail=str(e))

    # Anfrage-Log schreiben
    (logs_dir / f"{ts}_{qid}.json").write_text(stable_json_dumps({...}))
    return AskResponse(cad_metadata=cad, answers=list(answers))
\end{lstlisting}

\begin{infobox}[\texttt{asyncio.gather} -- Parallele Verarbeitung]
Jede Frage durchläuft vollständig isoliert: eigener Retrieve-Aufruf, eigener
LLM-Aufruf. \texttt{asyncio.gather} startet alle Coroutinen gleichzeitig und
wartet bis alle abgeschlossen sind. Bei $n$ Fragen beträgt die Gesamtdauer
annähernd die Zeit einer einzigen Frage (statt $n \times$).
\end{infobox}

\subsection{Upload-Endpunkt}

\begin{lstlisting}[style=pythonstyle, caption={PDF-Upload und Indizierung}]
@app.post("/upload")
async def upload_pdf(file: UploadFile = File(...)):
    # In temporaere Datei speichern
    dest = uploads_dir / f"{timestamp}_{uuid}.pdf"
    dest.write_bytes(await file.read())

    # Indizierung in separatem Thread (blockierender CPU-Code)
    try:
        info = await asyncio.to_thread(indexer.index_pdf, dest)
    except Exception as e:
        raise HTTPException(status_code=500, detail=f"Indexing failed: {e}")

    return info   # DocumentInfo: source_path, doc_hash, chunk_count
\end{lstlisting}

\texttt{asyncio.to\_thread} lagert den blockierenden Indizierungs-Code in einen
Thread-Pool aus, damit der Event-Loop für andere Anfragen frei bleibt.

\newpage
% ============================================================
\section{Frontend: \texttt{frontend/index.html}}
% ============================================================

Das gesamte Frontend ist eine einzige HTML-Datei ohne Build-Prozess.
Es verwendet Pico.css (klassenfreies CSS-Framework) und Vanilla JavaScript.

\subsection{Zustandsmodell}

\begin{lstlisting}[style=pythonstyle, caption={JavaScript-Zustandsmodell}]
const state = {
  cad: {},      // Aktuelle CAD-Metadaten (Objekt)
  maxQ: 6,      // Maximale Anzahl Fragefelder
  defaultQ: 3,  // Standard-Anzahl beim Start
};
\end{lstlisting}

\subsection{Fehlerbehandlung}

\begin{lstlisting}[style=pythonstyle, caption={Fehlerdarstellung (JavaScript)}]
function formatError(payload, fallback) {
  if (!payload) return fallback;
  const d = payload.detail ?? payload;  // FastAPI-Fehler haben "detail"-Feld
  if (typeof d === "string") return d;
  try { return JSON.stringify(d, null, 2); }
  catch (_) { return String(d); }
}
\end{lstlisting}

\subsection{Fragen-API-Aufruf}

\begin{lstlisting}[style=pythonstyle, caption={Frontend-Aufruf der /ask-API}]
el("ask").onclick = async () => {
  const questions = getQuestions();         // Alle nichtleeren Felder sammeln
  el("askStatus").textContent = "Laeuft...";

  const r = await fetch("/ask", {
    method: "POST",
    headers: {"Content-Type": "application/json"},
    body: JSON.stringify({
      questions,
      cad_metadata: state.cad,
      format: el("format").value
    })
  });

  const j = await r.json().catch(() => null);
  el("askStatus").textContent = r.ok ? "Fertig." : `Error: ${formatError(j, "...")}`;
  if (r.ok) renderOutput(j);
};
\end{lstlisting}

\newpage
% ============================================================
\section{Konfigurationsdatei: \texttt{config.yaml}}
% ============================================================

\begin{lstlisting}[style=yamlstyle, caption={Aktuelle Konfiguration (config.yaml)}]
domain:
  name: "Verzahnungen"
  schema_path: "schemas/gears.yaml"
  prompt_path: "prompts/answer_system_prompt.txt"

embedder:
  implementation: "bge_m3"
  model_name: "BAAI/bge-m3"
  device: "mps"            # Apple Silicon GPU
  use_sparse: false
  max_length: 8192

chunker:
  implementation: "recursive"
  min_chunk_size: 100      # Token-Untergrenze
  max_chunk_size: 512      # Token-Obergrenze
  overlap_sentences: 1     # Satz-Ueberlappung zwischen Chunks

metadata_extractor:
  model_name: "llama3.2:3b"
  ollama_url: "http://localhost:11434"
  max_retries: 2
  timeout_s: 30

vector_store:
  implementation: "qdrant"
  host: "localhost"
  port: 6333
  collection_name: "knowledge_base"

retriever:
  stage1_strict: true
  stage1_relax_on_empty: true  # Lockert Filter, wenn < 5 Treffer
  stage1_min_candidates: 5
  stage2_top_k: 5
  stage2_min_similarity: 0.65  # Mindest-Kosinusaehnlichkeit

answer_generator:
  model_name: "llama3.2:3b"
  default_format: "standard"
  max_tokens: 800
  temperature: 0.2             # Niedrig = deterministisch
\end{lstlisting}

% ============================================================
\section{Domänenschema: \texttt{schemas/gears.yaml}}
% ============================================================

Das Schema ist vollständig extern konfigurierbar. Ein Wechsel der Domäne
(z.\,B. Lager, Schweißverbindungen) erfordert nur eine neue YAML-Datei
und eine Pfadanpassung in \texttt{config.yaml}.

\begin{lstlisting}[style=yamlstyle, caption={Metadatenschema Verzahnungen}]
domain: "Verzahnungen"

fields:
  - name: verzahnungstyp
    type: string
    enum: ["Stirnrad", "Schraegverzahnung", "Kegelrad", ...]
    filter_type: exact       # Exakter Vergleich in Stage-1-Filter

  - name: modul
    type: number
    filter_type: range       # Bereichsvergleich
    range_fields: [modul_min, modul_max]

  - name: norm
    type: string
    filter_type: exact
    nullable: true           # Darf null sein

  - name: themenkategorie
    type: string
    enum: ["Geometrie", "Festigkeit", "Fertigung", ...]
    filter_type: set         # Mengenabgleich

  - name: anwendungsbereich
    type: string
    nullable: true
    # kein filter_type -> wird nicht im Stage-1-Filter verwendet
\end{lstlisting}

\newpage
% ============================================================
\section{Vollständiger Anfragepfad (Beispiel)}
% ============================================================

\begin{enumerate}[leftmargin=1.5cm, label=\textbf{Schritt \arabic*.}]
  \item \textbf{Frontend}: Nutzer stellt 2 Fragen, klickt ``Fragen stellen''
  \item \textbf{POST /ask}: Body mit \texttt{questions}, \texttt{cad\_metadata}, \texttt{format}
  \item \textbf{API}: Validierung durch Pydantic (\texttt{AskRequest})
  \item \textbf{asyncio.gather}: 2 parallele \texttt{\_answer\_one}-Coroutinen starten
  \item \textbf{Für jede Frage}:
    \begin{enumerate}[label=\alph*)]
      \item \texttt{retriever.retrieve(question, cad\_metadata)}
      \item Embedder: Frage → Dense-Vektor ($\mathbb{R}^{1024}$)
      \item Stage-1-Filter: CAD-Metadaten → Qdrant-Filter-Objekt
      \item Qdrant: \texttt{query\_points} mit Filter + Score-Threshold 0.65
      \item Falls $< 5$ Treffer: Filter mit Faktor 1.5, dann 2.0 lockern
      \item \texttt{answer\_generator.generate(question, chunks, cad, format)}
      \item Prompt-Template befüllen (DOMAIN, CAD, CHUNKS, QUESTION, FORMAT)
      \item Ollama: HTTP POST an \texttt{/api/generate} → LLM-Antworttext
      \item \texttt{Answer}-Dict zurückgeben (question, answer\_text, sources)
    \end{enumerate}
  \item \textbf{Logging}: Anfrage-Log in \texttt{logs/queries/\{ts\}\_\{uuid\}.json}
  \item \textbf{Response}: \texttt{AskResponse} mit \texttt{cad\_metadata} + \texttt{answers}
  \item \textbf{Frontend}: Antwort-Karten rendern mit Quellen-Akkordeon
\end{enumerate}

% ============================================================
\section{Erweiterungspunkte}
% ============================================================

Das System ist an definierten Punkten ohne Code-Eingriff erweiterbar:

\begin{table}[H]
\centering
\small
\begin{tabularx}{\textwidth}{lXl}
\toprule
\textbf{Erweiterung} & \textbf{Was zu tun ist} & \textbf{Aufwand} \\
\midrule
Neues Embedding-Modell & Neue Klasse mit \texttt{embed()}-Methode, Eintrag in \texttt{factory.py} & klein \\
Neues LLM & Neuer Client (z.\,B. OpenAI), Klasse mit \texttt{generate()} & klein \\
Neue Domäne & Neue \texttt{schemas/<domain>.yaml} + Config-Pfad ändern & minimal \\
Echter CAD-Adapter & Klasse mit \texttt{extract(file\_path)} -> dict & mittel \\
Reranker & 3. Stufe im Retriever (per Config aktivierbar) & mittel \\
Hybrid-Search & Sparse-Vektoren in Qdrant (bereits vorbereitet) & mittel \\
Multi-User & Auth-Middleware in FastAPI + User-Isolation in Qdrant & groß \\
\bottomrule
\end{tabularx}
\caption{Erweiterungspunkte und geschätzter Aufwand}
\end{table}

% ============================================================
\section{Abhängigkeiten}
% ============================================================

\begin{table}[H]
\centering
\small
\begin{tabularx}{\textwidth}{llX}
\toprule
\textbf{Bibliothek} & \textbf{Version} & \textbf{Verwendung} \\
\midrule
\texttt{fastapi} & -- & REST-API-Framework \\
\texttt{uvicorn} & -- & ASGI-Server \\
\texttt{pydantic} & v2 & Konfiguration, Datenvalidierung \\
\texttt{sentence-transformers} & -- & BAAI/bge-m3 Embedding \\
\texttt{qdrant-client} & 1.16+ & Vektordatenbank-Client \\
\texttt{PyMuPDF (fitz)} & -- & PDF-Textextraktion \\
\texttt{httpx} & -- & HTTP-Client für Ollama \\
\texttt{pyyaml} & -- & YAML-Konfiguration und Schema \\
\bottomrule
\end{tabularx}
\caption{Python-Abhängigkeiten}
\end{table}

Externe Dienste (müssen separat laufen):
\begin{itemize}
  \item \textbf{Qdrant}: Docker \texttt{docker-compose up} (Port 6333)
  \item \textbf{Ollama}: \texttt{ollama serve} + \texttt{ollama pull llama3.2:3b} (Port 11434)
\end{itemize}

\vfill
\begin{center}
  \small\color{accentgray}
  Automatisch generiert aus dem Quellcode des KI-Copiloten für Verzahnungswissen.
\end{center}

\end{document}
